% Vietnamese translation for OI Public Library
% Source-File: IOI中国国家候选队论文/2008/Day1/5.周梦宇《码之道——浅谈信息学竞赛中的编码与译码问题》/码之道.ppt
\documentclass[11pt]{article}
\usepackage{oipl-vn}

\newcommand{\Oh}{\mathcal{O}}

\title{Bản trình chiếu: Đạo của mã}
\author{Zhou Mengyu}
\date{2008}

\begin{document}
\maketitle

\origtitle{Slide companion for The Tao of Coding}
\sourcepath{IOI China candidate team papers / 2008 / Day 1 / Zhou Mengyu.ppt}

\section{Phạm vi bản dịch}

Tệp PowerPoint này là bản trình chiếu của bài \textit{Đạo của mã}; bản PPTX
cùng chủ đề đã được dịch. Bản ghi chú này tóm lược cách hiểu mã hóa và giải
mã trong thi tin học.

\section{Mã hóa trong OI}

Trong bối cảnh thi lập trình, mã hóa thường là biến một đối tượng tổ hợp thành
một số thứ tự, một chuỗi bit hoặc một biểu diễn ngắn. Giải mã là dựng lại đối
tượng từ biểu diễn ấy. Các bài kiểu này thường xoay quanh hoán vị, tập con,
chuỗi, đường đi hoặc cấu hình có ràng buộc.

\section{Đếm theo lớp}

Với thứ tự từ điển của hoán vị, sinh lần lượt bằng
\texttt{next\_permutation} là quá chậm. Cách tốt hơn là đếm số đối tượng trong
từng lớp tiền tố. Khi mã hóa, ta cộng số lớp nhỏ hơn tiền tố hiện tại; khi
giải mã, ta dùng thứ hạng để chọn lớp chứa đối tượng cần dựng.

\section{Công cụ}

Các slide gợi ý phối hợp đếm tổ hợp, quy hoạch động, tìm kiếm, sắp xếp và
tham lam. Điểm chung là mỗi bước mã hóa hoặc giải mã cần biết số đối tượng
trong một lớp, vì vậy bài toán thường trở thành bài toán đếm có điều kiện.

\end{document}
